\documentclass[11pt,reqno]{amsart}
\usepackage[T1]{fontenc}
\usepackage{lmodern}
\usepackage[margin=1in]{geometry}
\usepackage{microtype,amsmath,amssymb,mathtools}
\usepackage[hidelinks]{hyperref}
\hypersetup{pdftitle={Removing stationarity from logarithmic DLR uniqueness}}
\usepackage{enumitem}
\numberwithin{equation}{section}
\newtheorem{theorem}{Theorem}[section]
\newtheorem{proposition}[theorem]{Proposition}
\newtheorem{lemma}[theorem]{Lemma}
\newtheorem{corollary}[theorem]{Corollary}
\theoremstyle{definition}
\newtheorem{definition}[theorem]{Definition}
\theoremstyle{remark}
\newtheorem{remark}[theorem]{Remark}
\newcommand{\R}{\mathbb R}
\newcommand{\Z}{\mathbb Z}
\newcommand{\N}{\mathbb N}
\newcommand{\E}{\mathbb E}
\newcommand{\1}{\mathbf1}
\newcommand{\Sine}{\operatorname{Sine}}
\newcommand{\Conf}{\operatorname{Conf}}
\newcommand{\dd}{\,d}
\DeclareMathOperator{\supp}{supp}
\DeclareMathOperator{\diver}{div}
\allowdisplaybreaks[2]
\setlength{\emergencystretch}{2em}
\setlist[enumerate]{leftmargin=*,itemsep=3pt}
\title[Removing stationarity from logarithmic DLR uniqueness]{Removing stationarity from logarithmic DLR uniqueness}
\date{14 September 2026}

\begin{document}
\begin{abstract}
We prove that a canonical one-dimensional logarithmic DLR law at any
inverse temperature $\beta>0$ is translation-invariant whenever the sums
of squared unit-interval counts grow at most linearly, almost surely.
The hypothesis is pathwise; its random constant need not be integrable.
Finite spatial mean energy of one regularized electric field implies this
count condition. Consequently the finite-energy uniqueness theorem in the
accompanying manuscript extends to laws not assumed stationary, under
the same regularization convention. The argument uses compactly supported
translations with a logarithmic taper and an exact product identity for
two Radon--Nikodym derivatives. A Hellinger estimate avoids expectations
of the energy errors. We also record a separate nonstationary
$\beta=2$ consequence of deterministic conditional universality.
\end{abstract}
\maketitle

\noindent\textit{Scope.}
This is a supplement to the supplied manuscript
\emph{Uniqueness of the stationary logarithmic DLR state}, file
\texttt{logarithmic\_dlr\_uniqueness(2).tex}, denoted by [M] below.
The automatic-stationarity theorem and its proof are given in full.
Identification with $\Sine_\beta$ then invokes the stationary uniqueness
theorem of [M]; this supplement does not reprove that theorem.

\section{The nonstationary statements}

Write $\theta_a\gamma=\gamma-a$ and
$N_j(\gamma)=\gamma([j,j+1))$, $j\in\Z$.
We use the canonical logarithmic conditional densities of
Dereudre--Hardy--Lebl\'e--Ma\"ida~\cite{DHLM}: condition on both the
configuration outside a bounded interval $\Lambda$ and the number of
points inside it. On the ordered configuration simplex the conditional
density is positive almost everywhere and proportional to
$\exp(-\beta H_\Lambda)$, relative to Lebesgue measure.
All exterior energy differences use common symmetric spatial cutoffs.
In particular they exist as finite relative energies on the conditional
full-measure sets. We retain these Gibbs-kernel equations, without
including stationarity in the definition of a solution. This distinction
matters because \cite[Definition 3.1]{DHLM} packages stationarity and
energy assumptions together with the equations.

More explicitly, if $T$ is an increasing $C^1$ diffeomorphism equal to
the identity outside a compact subset of $\Lambda$, its energy change is
\begin{equation}\label{eq:move}
\begin{split}
\Delta H_T(\gamma)={}&
-\sum_{\{x,y\}\subset\gamma\cap\Lambda}
 \log\frac{|T(x)-T(y)|}{|x-y|}\\
&-\lim_{S\to\infty}
 \sum_{\substack{x\in\gamma\cap\Lambda\\
                   y\in\gamma\setminus\Lambda,\ |y|\le S}}
 \log\frac{|T(x)-y|}{|x-y|}.
\end{split}
\end{equation}
The first sum is over unordered distinct pairs. Normalizing constants
and reference configurations disappear in this difference.

\begin{definition}[A pathwise count condition]\label{def:temper}
A configuration is square-count tempered if
\begin{equation}\label{eq:temper}
M(\gamma):=\sup_{m\in\N,\ m\ge1}
\frac{1}{m+1}\sum_{|j|\le m}N_j(\gamma)^2<\infty.
\end{equation}
\end{definition}

\begin{theorem}[Automatic stationarity]\label{thm:auto}
Fix $\beta>0$. Let $P$ be a law on locally finite simple configurations
on $\R$ satisfying the canonical logarithmic Gibbs-kernel equations
described above. If \eqref{eq:temper} holds $P$-almost surely, then
\[
P\circ\theta_a^{-1}=P\qquad\text{for every }a\in\R.
\]
No moment assumption on $M(\gamma)$, no prescribed intensity, and no
number-rigidity assumption are required.
\end{theorem}

For the electric formulation, let $\rho_\eta$ be uniform probability
measure on the circle of radius $\eta$ in $\R^2$, and regard points of
$\gamma$ as points on the horizontal axis. Consider the pathwise condition
that there exist $0<\eta<1$ and $F\in L^2_{\mathrm{loc}}(\R^2;\R^2)$ with
\begin{align}
-\diver F&=2\pi(\gamma*\rho_\eta-dx\,\delta_0),\label{eq:div}\\
\limsup_{R\to\infty}\frac1{2R}
\int_{[-R,R]\times\R}|F|^2&<\infty.\label{eq:field}
\end{align}
Both $\eta$ and $F$ may depend on the configuration; a measurable
selection is unnecessary.

\begin{corollary}[Finite energy forces stationarity]\label{cor:energy-stationarity}
A law satisfying the canonical logarithmic Gibbs-kernel equations and
supported on configurations satisfying \eqref{eq:div}--\eqref{eq:field}
is stationary.
\end{corollary}

\begin{corollary}[Nonstationary finite-energy uniqueness]\label{cor:uniqueness}
Use the standard inner electric-field truncation and the renormalized
energy $W$ specified in [M, Section 7]. Let $P$, not assumed stationary,
be a simple point process satisfying the canonical logarithmic
Gibbs-kernel equations, and suppose $\E_P W<\infty$ in the same sense as
[M, Theorem 1.2]. Then, using that stationary theorem,
\[
P=\Sine_\beta.
\]
\end{corollary}

The proof of automatic stationarity is a localized-translation argument
of the type used for conservation of continuous symmetries; see
Georgii~\cite{Geo} and, for the two-dimensional plasma,
Lebl\'e~\cite[Section 8]{Leble}. Those results concern different settings.
We establish all estimates needed for the one-dimensional logarithmic
interaction directly below.

\section{Cutoffs with vanishing deformation cost}

For $f\in C_c^1(\R)$ define
\[
q_f(x,y)=\frac{f(x)-f(y)}{x-y}\quad(x\ne y),
\qquad q_f(x,x)=f'(x).
\]
The diagonal definition makes this a continuous divided difference.

\begin{lemma}[Deterministic cutoff estimate]\label{lem:cutoff}
Fix $A\ge2$. There exist even functions $f_n\in C_c^\infty(\R)$,
$0\le f_n\le1$, such that $f_n=1$ on $[-A,A]$,
$\supp f_n\subset[-R_n,R_n]$, $R_n=A2^n$, and
\begin{equation}\label{eq:lip}
\|f_n'\|_\infty\le\frac{C}{An}.
\end{equation}
For every configuration satisfying \eqref{eq:temper},
\begin{equation}\label{eq:cost}
\sum_{\{x,y\}\subset\gamma}q_{f_n}(x,y)^2
+\sum_{x\in\gamma}f_n'(x)^2
\le\frac{C_A M(\gamma)}{n}.
\end{equation}
The sums over pairs are unordered and exclude the diagonal.
\end{lemma}

\begin{proof}
Choose an even smooth function $\psi$, nonincreasing on $[0,\infty)$,
equal to one on $[-1,1]$ and zero outside $[-2,2]$. Put
\begin{equation}\label{eq:cutoff}
f_n(x)=\frac1n\sum_{k=0}^{n-1}\psi\!\left(\frac{x}{A2^k}\right).
\end{equation}
The derivative supports lie in the disjoint open annuli
$A2^k<|x|<A2^{k+1}$. In particular
\begin{align}
|f_n'(x)|&\le\frac{C}{n(A+|x|)},\label{eq:derivative}\\
|f_n(x)-f_n(y)|&\le\frac Cn
\left|\log\frac{A+|x|}{A+|y|}\right|.\label{eq:logdiff}
\end{align}
The second bound follows by integrating the first along the radial
variable and using evenness. The derivative vanishes outside
$[-R_n,R_n]$.

Let $I_j=[j,j+1]$ and set
\[
d_j=\sup_{I_j}|f_n'|,\qquad
b_j=\sum_{k\in\Z}\sup_{I_j\times I_k}|q_{f_n}|^2.
\]
We claim the row estimates
\begin{equation}\label{eq:row}
b_j\le\frac{C_A}{n^2}
\begin{cases}
(A+|j|)^{-1},& |j|\le4R_n,\\
R_n|j|^{-2},& |j|>4R_n.
\end{cases}
\end{equation}
Here and below constants may depend on the fixed core size $A$, but
never on $n$ or the configuration.

First, the elementary integral
\[
\int_0^\infty\frac{\log^2 t}{(1-t)^2}\dd t<\infty
\]
and the change of variable $t=(A+y)/(A+x)$ show, for $x\ge0$, that
\[
\int_0^\infty
\frac{\log^2((A+x)/(A+y))}{(x-y)^2}\dd y
\le\frac{C}{A+x}.
\]
For negative $y$, replace $y$ by $|y|$ and use
$|x-y|\ge|x-|y||$. Reflection covers $x<0$. Thus
\begin{equation}\label{eq:continuous-row}
\int_\R |q_{f_n}(x,y)|^2\dd y
\le\frac{C}{n^2(A+|x|)}.
\end{equation}

To pass to cell suprema, neighboring cells are bounded by the
mean-value theorem. For $|j-k|\ge2$, distances between the two cells
are comparable to $|j-k|$, and changing either endpoint to the left
endpoint of its cell changes $f_n$ by at most $d_j$ or $d_k$.
Comparing the remaining value at $k$ with its average over $I_k$ gives
\begin{equation}\label{eq:cell-row}
b_j\le C\int_\R |q_{f_n}(j,y)|^2\dd y
 +C\sum_{k\in\Z}\frac{d_j^2\1_{\{k=j\}}+d_k^2}
                         {(1+|j-k|)^2}.
\end{equation}
For clarity, the $d_j^2$ errors of all distant cells have already
been summed, using $\sum_{h\ne0}h^{-2}<\infty$.
Since $d_k\le C/[n(A+|k|)]$, splitting the convolution into
$|k|\ge|j|/2$ and $|k|<|j|/2$ bounds the second term in
\eqref{eq:cell-row} by $C_A/[n^2(A+|j|)^2]$.
Together with \eqref{eq:continuous-row}, this proves the first row bound
(in fact that bound holds for every $j$).

For the summable far-tail improvement, observe that
\begin{equation}\label{eq:l2cells}
\sum_{k\in\Z}\sup_{I_k}f_n^2\le\frac{C_A R_n}{n^2}.
\end{equation}
Indeed, on the annulus $A2^\ell\le|x|\le A2^{\ell+1}$,
$f_n(x)\le(n-\ell)/n$. Summing its squared upper bound times the
number of cells gives
$C_A n^{-2}\sum_{\ell=0}^{n-1}2^\ell(n-\ell)^2
\le C_A R_n/n^2$; the fixed central interval obeys the same bound
after increasing $C_A$. Cells meeting annulus endpoints are covered
by enlarging the annuli by one unit. If $|j|>4R_n$, $f_n=0$ on $I_j$
and every contributing $I_k$ has distance comparable to $|j|$ from
$I_j$. Equation~\eqref{eq:l2cells} proves the second row bound.

Put $a_j=N_j(\gamma)^2$. Since
$N_jN_k\le(a_j+a_k)/2$, symmetry of the cell kernel gives
\begin{equation}\label{eq:weighted-row}
\sum_{\{x,y\}\subset\gamma}q_{f_n}(x,y)^2
\le\frac12\sum_j a_j b_j.
\end{equation}
The harmless diagonal overcount only increases the right side.
By \eqref{eq:temper}, the sum of the $a_j$ in a dyadic shell of
radius $s\ge1$ is at most $C M(\gamma)s$.
Each shell inside $4R_n$ therefore contributes at most
$C_A M(\gamma)/n^2$ to \eqref{eq:weighted-row}, and there are
$n+O_A(1)$ such shells. Beyond $4R_n$, the shell of radius $2^\ell R_n$
contributes at most $C_A M(\gamma)2^{-\ell}/n^2$.
This proves the pair part of \eqref{eq:cost}.

Finally $N_j\le N_j^2$, because counts are nonnegative integers.
The derivative bound on each annulus gives
\[
\sum_{x\in\gamma}f_n'(x)^2
\le\sum_j a_jd_j^2
\le\frac{C_A M(\gamma)}{n^2},
\]
by the same dyadic decomposition, now with summable shell contributions.
This completes the proof.
\end{proof}

\begin{remark}
The essential feature is the critical divided-difference kernel
$|x-y|^{-2}$. Distributing the change over $n$ dyadic scales costs
$O(n^{-1})$. One may also view \eqref{eq:cutoff} as the trace of radial
cutoffs in $\R^2$ with Dirichlet energy $O(n^{-1})$, that is, vanishing
homogeneous $H^{1/2}$ cutoff energy for a fixed interval. The cell estimates above are
needed to accommodate arbitrary locally finite configurations.
\end{remark}

\section{DLR changes of variables and the Hellinger estimate}

\begin{proof}[Proof of Theorem~\ref{thm:auto}]
Fix a translation $a\in\R$ and a bounded measurable local observable
$G$, determined by the configuration in a bounded interval $B$.
Choose $A\ge2$ with $B\subset[-A,A]$, and use Lemma~\ref{lem:cutoff}.
For all sufficiently large $n$, define
\[
T_{n,\pm}(x)=x\pm af_n(x).
\]
By \eqref{eq:lip} we may assume
$|a|\|f_n'\|_\infty\le1/2$. These maps are increasing smooth
diffeomorphisms, equal to the identity outside $[-R_n,R_n]$.
They act on configurations by moving every point.
Let
\[
Q_{n,\pm}=(T_{n,\pm}^{-1})_*P.
\]
We deliberately use the inverse maps in this pushforward.

Choose a bounded interval $\Lambda_n$ containing $[-R_n,R_n]$ in
its interior. Each $T_{n,\pm}$ maps $\Lambda_n$ onto itself, fixes
its exterior, and preserves its number of points. On every conditional
fiber of $(\gamma_{\Lambda_n^c},N_{\Lambda_n})$, the canonical DLR
density and ordinary change of variables therefore give
\begin{equation}\label{eq:RN}
w_{n,\pm}(\gamma):=\frac{dQ_{n,\pm}}{dP}(\gamma)
=\exp\{-\beta\Delta H_{T_{n,\pm}}(\gamma)\}
 \prod_{x\in\gamma}(1\pm af_n'(x)).
\end{equation}
The product is finite. Positivity of the conditional densities and
the diffeomorphism property give equivalence of the measures and
$\E_Pw_{n,\pm}=1$. Conditional exceptional Lebesgue-null sets remain
null under these diffeomorphisms.

Put $q_n=q_{f_n}$. For every distinct pair $x,y$,
\[
\frac{|T_{n,+}(x)-T_{n,+}(y)|}{|x-y|}
\frac{|T_{n,-}(x)-T_{n,-}(y)|}{|x-y|}
=1-a^2q_n(x,y)^2.
\]
Applying this identity to the interior pairs and then to the exterior
pairs at a common finite cutoff yields
\begin{equation}\label{eq:product}
w_{n,+}w_{n,-}=e^{-\mathcal A_n},
\end{equation}
where
\begin{equation}\label{eq:An}
\mathcal A_n(\gamma)=
-\beta\sum_{\{x,y\}\subset\gamma}\log(1-a^2q_n(x,y)^2)
-\sum_{x\in\gamma}\log(1-a^2f_n'(x)^2).
\end{equation}
There is no uncontrolled linear exterior-force term: it cancels at
the common cutoff before the limit is taken. The individual relative
energies exist by the specification. The paired infinite sum is
absolutely convergent by \eqref{eq:cost} and the bound below, so passing
to the limit preserves \eqref{eq:product}.

The mean-value theorem gives $|a q_n(x,y)|\le1/2$. Using
$0\le-\log(1-z)\le2z$ for $0\le z\le1/4$, we obtain, $P$-almost surely,
\begin{equation}\label{eq:An-small}
0\le\mathcal A_n(\gamma)
\le2a^2\left(\beta\sum_{\{x,y\}\subset\gamma}q_n(x,y)^2
             +\sum_{x\in\gamma}f_n'(x)^2\right)
\le\frac{C_{A,a,\beta}M(\gamma)}n\longrightarrow0
\end{equation}
For $Q\ll P$ use the convention
$\mathsf H^2(Q,P)=\int(\sqrt{dQ/dP}-1)^2\dd P$.
The arithmetic--geometric mean inequality and \eqref{eq:product} give
\begin{align}
\mathsf H^2(Q_{n,+},P)+\mathsf H^2(Q_{n,-},P)
&=4-2\E_P(\sqrt{w_{n,+}}+\sqrt{w_{n,-}})\notag\\
&\le4\E_P[1-e^{-\mathcal A_n/4}]\longrightarrow0.\label{eq:hellinger}
\end{align}
The last step is bounded convergence. In particular we have not
assumed $\E_PM<\infty$, $\E_P\mathcal A_n<\infty$, or integrability
of individual energy changes. Cauchy--Schwarz gives
\[
\|Q-P\|_{\mathrm{TV}}
=\frac12\int|dQ/dP-1|\dd P\le\mathsf H(Q,P),
\]
so $Q_{n,+}\to P$ in total variation.

Since $T_{n,+}(x)=x+a$ on $B$, the restriction of
$T_{n,+}^{-1}\gamma$ to $B$ is exactly the restriction of
$\theta_a\gamma$ to $B$. The assertion includes particles entering the
interval: $T_{n,+}(B)=B+a$, and the map is one-to-one. Hence
\[
\E_{Q_{n,+}}G=\E_P[G\circ\theta_a].
\]
The left side tends to $\E_PG$. Equality for every bounded local
observable identifies the two probability measures, and $a$ was
arbitrary. This proves stationarity.
\end{proof}

\section{Why finite electric energy supplies the count condition}

\begin{lemma}\label{lem:energy-count}
Every locally finite configuration satisfying
\eqref{eq:div}--\eqref{eq:field} satisfies \eqref{eq:temper}.
\end{lemma}

\begin{proof}
Choose $\chi\in C_c^\infty((-2,3)\times(-2,2))$ with
$0\le\chi\le1$ and $\chi=1$ on $[-1,2]\times[-1,1]$.
Write $\chi_j(x,y)=\chi(x-j,y)$ and
$S_j=[j-2,j+3]\times[-2,2]$. Since $0<\eta<1$, every smearing circle
whose center lies in $[j,j+1)\times\{0\}$ is contained in
$\{\chi_j=1\}$. Positivity and the divergence identity give
\begin{align*}
N_j&\le\int\chi_j\dd(\gamma*\rho_\eta)\\
&=\int_\R\chi_j(x,0)\dd x
  +\frac1{2\pi}\int_{\R^2}F\cdot\nabla\chi_j\\
&\le5+C\left(\int_{S_j}|F|^2\right)^{1/2}.
\end{align*}
After squaring and summing, the boxes $S_j$ have bounded overlap, so
\begin{equation}\label{eq:count-field}
\sum_{|j|\le m}N_j^2
\le C(m+1)+C\int_{[-m-3,m+3]\times\R}|F|^2.
\end{equation}
The right side is $O_\gamma(m)$ by \eqref{eq:field}; the finitely
many remaining small $m$ have finite counts. This proves the lemma.
\end{proof}

\begin{proof}[Proof of Corollaries~\ref{cor:energy-stationarity}
and \ref{cor:uniqueness}]
Lemma~\ref{lem:energy-count} and Theorem~\ref{thm:auto} prove the
first corollary. For the second, recall the convention of [M]:
\[
f_\eta(z)=\1_{\{|z|\le\eta\}}\log(|z|/\eta),\qquad
E_\eta=E+\sum_{p\in\gamma}\nabla f_\eta(\,\cdot-(p,0)).
\]
If $-\diver E=2\pi(\gamma-dx\,\delta_0)$, then $F=E_\eta$
satisfies \eqref{eq:div}.

Finiteness of the renormalized energy, defined as the infimum over
compatible fields of their small-truncation renormalized spatial mean
energy, provides a compatible field for which this functional is
finite. The infimum need not be attained. For at least one sufficiently
small fixed $\eta$, its spatial mean of $|E_\eta|^2$ is finite:
the self-energy subtraction at fixed $\eta$ is a finite constant.
This is exactly the fixed-truncation implication used in
[M, Lemma 7.2]. It is pathwise and requires no stationarization.
Thus \eqref{eq:div}--\eqref{eq:field} hold $P$-almost surely.

The preceding corollary proves that $P$ is stationary. Its assumed
$\E_PW<\infty$ and its canonical DLR equations are unchanged.
It now satisfies all hypotheses of [M, Theorem 1.2], and hence
$P=\Sine_\beta$.
\end{proof}

\begin{remark}[The precise energy assumption]
The inner truncation indicator is essential; we use the convention
already specified in [M], rather than the outer indicator printed in
the supplied version of \cite[Definition 2.20]{DHLM}.
For an alternative notion of ``finite energy'', the implication
\eqref{eq:div}--\eqref{eq:field} must be checked. It is not enough
merely to produce a stationary spatial average of the candidate law.
\end{remark}

\section{What changes in the accompanying manuscript}

The argument supplies a preliminary stationarity theorem. The
gap-entropy comparison, random particle-ended blocks, Palm identities,
and stationary electric lift in [M] can subsequently be used as
written. They are applied only after stationarity has been proved.

In particular, the finite-energy theorem can be restated as follows:
\begin{quote}
Fix $\beta>0$, with the inner-truncation energy convention of [M].
Every simple point-process law satisfying the canonical logarithmic
Gibbs-kernel equations and $\E_PW<\infty$ equals $\Sine_\beta$.
\end{quote}
Its new proof first invokes Corollary~\ref{cor:energy-stationarity}
and then the original stationary theorem. The formulation must refer
to the Gibbs-kernel equations themselves, rather than invoking a
definition which already assumes stationarity.

For the discrepancy theorem, a similarly useful sufficient condition
is \eqref{eq:temper} together with the three count bounds in [M],
with their variance function replaced by
$d_P(t)=\E_P(N_{(0,t]}-t)^2$. Theorem~\ref{thm:auto}
first supplies stationarity; the linear centered bound then forces
intensity one, and the original comparison theorem applies.

The reason this additional argument is needed is structural. Averages
of many gap observables cannot distinguish a law from another law
with different behavior in a fixed window but the same spatial
averages. Stationarity is used when [M] identifies the limiting
averages with a Palm gap distribution. The present proof directly
controls translation of a fixed window and therefore resolves that
missing step. No assertion about a local defect being a DLR solution
is needed for this observation.

\section{An independent nonstationary result at \texorpdfstring{$\beta=2$}{beta=2}}

\begin{proposition}\label{prop:beta2}
Let $P$ satisfy the canonical logarithmic Gibbs-kernel equations at
$\beta=2$. Assume that, almost surely, its points admit an increasing
two-sided enumeration $(p_k)_{k\in\Z}$ with
\[
p_k/k\longrightarrow1\quad(k\to\pm\infty),\qquad
\lim_{S\to\infty}\sum_{\substack{p\in\gamma\\0<|p|\le S}}\frac1p
\text{ exists}.
\]
Then $P=\Sine_2$, without stationarity or an energy assumption.
\end{proposition}

\begin{proof}
Fix a configuration satisfying these assumptions. For the interval
$[-R,R]$ let $N_R=\gamma([-R,R])$. Its canonical conditional law is
the $N_R$-particle orthogonal-polynomial ensemble with weight
\[
w_R(x)=\1_{[-R,R]}(x)
\lim_{S\to\infty}\prod_{\substack{p\in\gamma\\R<|p|\le S}}
\left(1-\frac{x}{p}\right)^2.
\]
Kuijlaars--Mi\~na-D\'iaz~\cite[Assumption 1.2 and Theorem 1.3]{KMD}
prove that the associated kernels $K_R$ converge locally uniformly to
\[
K_{\sin}(x,y)=\frac{\sin\pi(x-y)}{\pi(x-y)},\qquad K_{\sin}(x,x)=1.
\]
Their particle number is exactly the actual count $N_R$ above.

For $h\in C_c^+(\R)$ supported in a compact interval $B$, the
conditional Laplace functional has the Fredholm series
\[
L_R(\gamma)=\sum_{k=0}^{N_R}\frac{(-1)^k}{k!}
\int_{B^k}\prod_{i=1}^k(1-e^{-h(x_i)})
\det[K_R(x_i,x_j)]_{i,j=1}^k\dd\mathbf x.
\]
For each fixed admissible configuration, eventually
$\sup_{x\in B}K_R(x,x)\le2$. Positive semidefiniteness and
Hadamard's inequality bound the absolute $k$th summand by
$(2|B|)^k/k!$. Local uniform convergence and dominated convergence
in the series therefore yield
$L_R(\gamma)\to\E_{\Sine_2}e^{-\langle h,\gamma\rangle}$.
Finally DLR gives
\[
\E_Pe^{-\langle h,\gamma\rangle}=\E_PL_R(\gamma),
\]
and $0\le L_R\le1$ allows bounded convergence over $P$.
Equality of all such Laplace functionals proves the proposition.
The random value of $N_R$ causes no difficulty: the deterministic
theorem is applied separately to each full configuration.
\end{proof}

This last result provides a separate check on the nonstationary
conclusion at $\beta=2$. The all-$\beta$ conclusion above instead
comes from the explicit automatic-stationarity argument and the
stationary theorem in [M].

\begin{thebibliography}{DHLM}
\bibitem[DHLM]{DHLM}
D. Dereudre, A. Hardy, T. Lebl\'e and M. Ma\"ida,
\emph{DLR equations and rigidity for the Sine-beta process},
arXiv:1809.03989v2 (2019).
\href{https://arxiv.org/abs/1809.03989v2}{arXiv:1809.03989v2}.

\bibitem[Geo]{Geo}
H.-O. Georgii,
\emph{Translation invariance and continuous symmetries in
two-dimensional continuum systems},
in \emph{Mathematical Results in Statistical Mechanics},
World Scientific (1999), 53--69.
\href{https://www.mathematik.uni-muenchen.de/~georgii/papers/D2sym.pdf}{Author's manuscript}.

\bibitem[KMD]{KMD}
A. B. J. Kuijlaars and E. Mi\~na-D\'iaz,
\emph{Universality for conditional measures of the sine point process},
J. Approx. Theory \textbf{243} (2019), 1--24.
\href{https://arxiv.org/abs/1703.02349v2}{arXiv:1703.02349v2},
especially Assumption 1.2 and Theorem 1.3.

\bibitem[Leb]{Leble}
T. Lebl\'e,
\emph{DLR equations, number-rigidity and translation-invariance
for infinite-volume limit points of the 2DOCP},
\href{https://arxiv.org/abs/2410.04958}{arXiv:2410.04958}, Section 8.
\end{thebibliography}
\end{document}
